\documentstyle[11pt]{article}
\setlength{\textheight}{8.8in}
\setlength{\columnsep}{2.0pc}
\setlength{\textwidth}{6.8in}
\setlength{\footheight}{0.0in}
\setlength{\topmargin}{0.25in}
\setlength{\headheight}{0.0in}
\setlength{\headsep}{0.0in}
\setlength{\oddsidemargin}{-.19in}
\setlength{\parindent}{1pc}

\title{An Hierarchy Description of the Blum's Complexity
of Languages}
\author{Wang Yongge\\
Nankai Institute of Mathematics\\
Nankai University\\
300071, Tianjin, P.R.China
}
\date{ }
\begin{document}
\date{}
\maketitle
\begin{abstract}
In this paper, we'll introduce a kind 
of $O(F)$-$LOOP$ operator,
which will lead to a hierarchy of any recursive functions:
$O(F)=\cup^\infty_{n=0}O(F)_n$. We've proved that this
hierarchy corresponds to the speed-up theorem, 
that is, given any 
$r(x)\in O(F)_i$, there is a language having $O(F)_{i+1}$ 
complexity which can be speeded up by $r(x)$, 
and also there is a
$r(x)$ in $O(F)_i$, so that any language having $O(F)_i$
complexity cannot be speeded up by $r(x)$. 
Through this, we can
grasp the soul of Speed-up theorem in a more higher level.
\end{abstract}

\section{Introduction}

In paper [6], Grzegorczyk gives a 
method to get the hierarchy
of the primitive recursive functions, and in [4], there are
many other methods to get the hierarchy of the primitive 
recursive functions. It has turned out that all 
these methods are 
equivalent. But it still remains a big question 
that how we can 
generalize these methods to the functions beyond 
the primitive 
recursive functions. In this paper, we'll introduce 
a kind of 
$O(F)$-$LOOP$ operator, which is a generalization 
of the methods
mentioned above and by which we can get a 
hierarchy of any set of 
recursive functions. The Grzegorczyk's
hierarchy and the hierarchy derived from the 
$O(F)$-$LOOP$ operator
have very good property, i.e. they correspond to 
the Speed-up 
theorem, from this the abstract Speed-up hierarchy 
can be defined as follows:
Given ${\cal F}=\cup^\infty_{n=0}{\cal F}_n$, if it 
satisfies some 
certain conditions, then the the sequence: ${\cal F}_1,{\cal F}_2,
\cdots$ is called a Speed-up hierarchy of 
${\cal F}$. At the end of this paper,
the following theorem is proved: for any 
$r(x)\in{\cal F}_i$,
there exists a language having ${\cal F}_{i+1}$ 
complexity, which 
can be speeded up by $r(x)$; and there exists 
another $r(x)\in{\cal F}_i$, such that any language having ${\cal F}_i$ 
complexity cannot 
be speeded up by $r(x)$. Because of this theorem, 
we can
understand the Speed-up theorem in a new sense, 
Speed-up theorem describes a
hierarchic speed-up phenomenon: for any speed-up 
factor $r(x)$, the language
$L$ that can be speeded up by $r(x)$ has at 
least one-level higher
complexity than $r(x)$, and the new algorithm 
got from the speed-up
still has the complexity of the same level as the old
algorithm, so we can really say that $r(x)$ does 
not speed up the
algorithm in essential. For example, the language 
that can be 
speeded up by polynomial has exponential(or higher) 
complexity and
the new algorithm got from the speed-up still has 
exponential
complexity.

\section{Program Language $O(F)$-$LOOP$}

Study in the following sections is based on 
a kind of $O(F)$-$LOOP$
language, so we first give the description 
of the language. In 
this language, some alphabets are used to 
denote variables.
Especially, let
$$X_1,X_2,\cdots\ \mbox{ and }\ Y$$
denote input and output variables, respectively; and let
$$Z_1,Z_2,\cdots$$
denote local variables.

The language $O(F)$-$LOOP$ has the following six types of
instructions:


1. $V\gets 0$ \ \ \ \ \ \ \ \ \ \ \ \ \ \ \ \ \ \ \ \ \ \ \ \ \ \ \ zero,

2. $V\gets V+1$ \ \ \ \ \ \ \ \ \ \ \ \ \ \ \ \ \ \ \ \ \ increment,

3. $V\gets V'$ \ \ \ \ \ \ \ \ \ \ \ \ \ \ \ \ \ \ \ \ \ \ \ \ \ assignment,

4. $LOOP \ V$ \ \ \ \ \ \ \ \ \ \ \ \ \ \ \ \ \ \ \ \ \ \ \ loop,

5. $END$ \ \ \ \ \ \ \ \ \ \ \ \ \ \ \ \ \ \ \ \ \ \ \ \ \ \ \ \ end,

6. $V\gets 0(f(V_1,\cdots,V_n))$ \ \ \ \ \ \ \ oracle;

\noindent
where $f\in F$ is a finite set of functions and
$V,V',V_1,\cdots,V_n$ are any variables.

Definition 1.1. $F$ is a given set of functions, a function
$g(x_1,\cdots,x_n)$ can be computed by a $O(F)$-$LOOP$ program $P$
iff for any input $X_1=c_1,X_2=c_2,\cdots,X_n=c_n$, the output 
$Y$ of an execution of $P$ is $g(c_1,\cdots,c_n)$.

Example: Given program $P$

$$\begin{array}{ll}
Z_1\longleftarrow X_1+1\\
Z_2\longleftarrow Z_1+1\\
\ \ \ LOOP\ \ Z_2\\
Z_2\longleftarrow Z_2+1\\
\ \ \ \mbox{END}\\
Z_3\longleftarrow O(f(Z_2))\\
Y\longleftarrow Z_3\end{array}$$
where $f(x)\in F$.

For input $x_1=1,x_2=2$, when $P$ runs to the third instruction
$LOOP\ \  Z_2$, the value of $Z_2$ is 3, so the loop body composed of
the third, fourth and fifth instructions should be executed 3 times. When
we are through with the execution of $P$, it will output 
$y=f(6)$. In fact, this program can compute the function
$y=f(2x_1+4)$.

Renmark: the fourth instruction in loop body will change the
value of $Z_2$, but this will not affect the number of loop 
iterations in the computation any more.

Definition 1.2. $O(F)$-$LOOP_n$ is the class of $O(F)$-$LOOP$
programs with $LOOP$---$END$ pairs nested to a depth of at most $n$.
$O(F)_n$ is the class of functions computable by programs in
$O(F)$-$LOOP_n$ and $O(F)$ is the class of functions computable by
$O(F)$-$LOOP$ programs.

Obviously, $\cup^\infty_{n=0}O(F)$-$LOOP_n$ is the class of all
$O(F)$-$LOOP$ programs and $O(F)=\cup^\infty_{n=0}O(F)_n$.
\vfill\pagebreak

\section{$E$-Ackermann Function}

For the reason of convenience, in the following sections, we
assume that: $F=\{f_1,\cdots,f_m\}$ is the set of 
functions that
satisfies the following properties (1) to (3)

(1). $f_i(x_1,\cdots,x_t)$ is a total function, 
$1\le i\le m$;

(2). $f_1(x)=x+2, f_2(x,y)=\max\{x,y\}$;

(3). for any component $x_i, f_j(x_1,\cdots,x_n)$ is
a nondecreasing function of $x_i$, $1\le i\le n, 2\le j\le m$.


Definition 2.1. the $E$-Ackermann function 
$A_F(m,n)$ based on 
$F$ is defined as follows:
$$\begin{array}{l}
A_F(0,0)=1,\\
A_F(0,1)=2,\\
A_F(0,n)=\max\{f_1(n),f_2(n,n),\cdots,f_m(n,\cdots,n)\}\ \ (n\ge 3),\\
A_F(m+1,n+1)=A_F(m,A(m+1,n));\end{array}$$
and the functions in sequence $\{u_n(x)|u_n(x)=A_F(n,x)\}$ are called
the hierarchy functions derived from $A_F(m,n)$.


We can easily show
$$\begin{array}{l} 
u_0(0)=1,\ \ u_0(1)=2,\\
u_0(x)=\max\{x+2,x,f_3(x,\cdots,x),f_m(x,\cdots,x)\}\ \ (x\ge 3);\\
u_{n+1}(x)=u^{(x)}_n(1)\end{array}$$
where $u^{(x)}_n(1)$ stands for $\underbrace{u_n\cdots u_n}_{x\
\mbox{ times }}(1)$.

Now we give some properties of $\{u_n(x)\}$ that will be used in 
the following sections.

Proposition 2.1. $u_{n+1}(x+1)=u_n(u_{n+1}(x))$.

Proposition 2.2. $u^{(k)}_0(x)\ge k$.

Proposition 2.3. $u_n(x)>x$.

Proposition 2.4. $u_n(x+1)>u_n(x)$.

Proposition 2.5. $u_{n+1}(x)\ge u_n(x)$.

Proposition 2.6. $u^{(k+1)}_n(x)>u^{(k)}_n(x)$.

Proposition 2.7. $u^{(k+1)}_n(x)\ge 2\cdot u^{(k)}_n(x),\ \ (n\ge 
1)$.

Proposition 2.8. $u^{(k+1)}_n(x)\ge u^{(k)}_n(x)+x$.

Proposition 2.9. $u^{(k)}_1(x)\ge 2^k\cdot x$.

The proofs of all the Propositions mentioned above are completely
like that of Ackermann function's properties, the reader is
refered to [1].

Proposition 2.10. For any $n,k,u_{n+1}(x)>u^{(k)}_n(x)$ a.e..

{\it Proof}. For $n=0$, we have
$$u_0(x)>x+2,\ \ u_1(m)=u^{(m)}_0(1)\ge 1+2m,$$
so for any $k$, when $x>2k-1$,
$$u^{(x-k)}_0(1)\ge 1+2(x-k)=1+2x-2k>x$$
i.e. for $x>2k-1$
$$u_1(x)=u^{(k)}_0(u^{(x-k)}_0(1))>u^{(k)}_0(x)$$

For $n\ne 0$, we proceed by induction on $k$.
For $k=0$, we have
using Proposition 2.3,
$$u_{n+1}(x)>x=u^{(0)}_n(x)$$

Suppose that $u_{n+1}(x)>u^{(k)}_n(x)$ a.e., then
$$\begin{array}{ll}  
&u^{(k+1)}_n(x)<u^{(k+1)}_n(2x-4)\ \ a.e.\\
\le &u^{(k+1)}_n(u_1(x-2))\ \ \ \ \ \ \mbox{ (Proposition 2.9)}\\
\le &u^{(k+1)}_n(u_n(x-2))\\
= &u^{(2)}_n(u^{(k)}_n(x-2))\\
\le &u^{(2)}_n(u_{n+1}(x-2))\ \ \ \ a.e.\\
=&u^{(2)}_n(u^{(x-2)}_n(1))\\
=&u_{n+1}(x)\end{array}$$


Proposition 2.11. $u_n(x)>x^2$ a.e. when $n>2$.

{\it Proof}. $$\begin{array}{l}  
u_1(x)\ge 2x,\\
u_2(x)\ge 2^x>x^2\ \mbox{ a.e.},\\
u_n(x)>u_2(x)>x^2\ \mbox{ a.e.}\ (n>2).\end{array}$$

\section{$O(F)$-$LOOP$\ \ Functions}

Definition 3.1. If $P$ is a $O(F)$-$LOOP$ program with input
variables $X_1,X_2,\cdots$, $X_n$, then the time, space and time-space
complexity of $P$ are defined as follows:
$T_P(x_1,\cdots,x_n)$ is the time complexity of $P$ defined as 
the total number of executions of zero, increment, assignment and 
oracle statements in the given computation. $S_P(x_1,\cdots,x_n)$ 
is the space complexity of $P$ defined as the maximal value of  
all the input variables, local variables and output variable in 
the given computation. And 
$ST_P(x_1,\cdots,x_n)=\max\{S_P(x_1,\cdots,x_n),T_P(x_1,\cdots,x_n)\}$ 
is the time-space complexity of $P$.


Theorem 3.1. If $P\in O(F)$-$LOOP_n$ then $ST_P(x_1,\cdots,x_n)\in
O(F)_n$.

{\it Proof}. To compute $ST_P$ for a given program $P$, we simply
modify $P$ by inserting counters $T$ and $S$, where $T$ is 
incremented each time a zero, increment, assignment or an oracle 
statement is executed in $P$, and $S$ is incremented according to 
the value of variables in the execution of $P$. This is 
accomplished by placing instructions ``$T\gets T+1$'' and 
``$S\gets O(\max(S,V))$'' immediately after each zero, increment, 
assignment and oracle statement in $P$, where $V$ is the variable 
whose value is changed in the statement after which ``$S\gets 
O(\max(S,V))$'' is inserted. Lastly, place an instruction 
``$Y\gets O(\max(T,S))$'' at the end of $P$. Thus, at the end of 
the computation, $S$, $T$ and $Y$ contains the value of $S_P,T_P$
and $ST_P$. This modified program obviously has  the same depth 
of nesting as $P$.



Theorem 3.2. If $P\in O(F)$-$LOOP_n$, then there is a $k$ such that
$$ST_P(x_1,\cdots,x_m)\le u^{(k)}_n(\max(x_1,\cdots,x_m)).$$

{\it Proof}. Let $v$ denote $\max(x_1,\cdots,x_m)$.
By induction on $n$. For $n=0$, there are no loops in $P$,
suppose that there are $k$ statements in $P$, so
$$ST_P(x_1,\cdots,x_m)\le u^{(k)}_{n-1}(v)\le u^{(k)}_n(v)$$
where we have used Proposition 2.3.

Now assume that the result is known for $n-1$ and let 
$P\in O(F)$-$LOOP_n$, we consider three cases:

{\it Case} 1. $P\in O(F)$-$LOOP_{n-1}$, using Proposition 2.6. and
induction hypothesis.
$$ST_P(x_1,\cdots,x_m)\le u^{(k)}_{n-1}(v)\le u^{(k)}_n(v)$$

{\it Case} 2. $P$ has the form
$$LOOP\ \ V$$
$$Q$$
$$END\ \ \ \ \ \ $$
where $Q\in O(F)$-$LOOP_{n-1}$. By induction hypothesis, there is a
$k$ such that
$$S_Q(x_1,\cdots,x_m)\le ST_Q(x_1,\cdots,x_m)\le u^{(k)}_{n-
1}(v)$$
$$T_Q(x_1,\cdots,x_m)\le ST_Q(x_1,\cdots,x_m)\le u^{(k)}_{n-
1}(v)$$

Hence the possible maximal value of all variables after the first 
running of loop body $Q$ is
$$u^{(k)}_{n-1}(v)$$
and after the second running of loop body $Q$, the possible 
maximal value of all variables is
$$u^{(2k)}_{n-1}(v)$$
proceeding in the same way, we have
$$\begin{array}{l}  
S_P(x_1,\cdots,x_m)\\
\le u^{(v\cdot k)}_{n-1}(v)\\
\le u^{(vk+1)}_{n-1}(v)\\
\le u^{(v\cdot(k+1))}_{n-1}(1)\\
\le u^{(v\cdot 2^{k+1})}_{n-1}(1)\\
\le u^{(u^{(k+1)}_n(v))}_{n-1}(1)\\
\le u^{(u^{(k+1)}(v))}_{n-1}(1)\\
=u_n(u^{(k+1)}_n(v))\\
=u^{(k+2)}_n(v),\end{array}$$
using Proposition 2.8. repeatedly, we have
$$\begin{array}{l}  
T_P(x_1,\cdots,x_m)\\
\le u^{(k)}_{n-1}(v)+u^{(2k)}_{n-1}(v)+\cdots+u^{(v\cdot k)}_{n-
1}(v)\\
\le u^{(k+1+k)}_{n-1}(v)+u^{(3k)}_{n-1}(v)+\cdots+u^{(v\cdot 
k)}_{n-1}(v)\\
\cdots\cdots\\
\le u^{(vk+1)}_{n-1}(v)\\
\cdots\cdots\\
\le u^{(k+2)}_n(v).\end{array}$$
hence
$$ST_P(x_1,\cdots,x_m)\le u^{(k+2)}_n(v)$$

{\it Case} 3. For the general case, $P$ has the form
$$\begin{array}{l} 
\ \ Q_0\\
LOOP\ \ V_1\\
\ \ P_1\\
END\\
\ \ \cdots\\
LOOP\ \ V_r\\
\ \ P_r\\
END\\
\ \ Q_r\end{array}$$
where $P_i,Q_i\in O(F)$-$LOOP_{n-1}$.


Using Proposition 2.8. repeatedly, we have
$$\begin{array}{l}  
ST_P(x_1,\cdots,x_m)\\
\le 
u^{(k_1)}_n(v)+u^{(k_2)}_n(u^{(k_1)}_n(v))+\cdots+u^{(k_s)}_n(u^{(k_{s-1})}_n
(\cdots u^{(k_1)}_n(v)\cdots))\\
\le\cdots\le u^{(k_s+\cdots+k_1+s-1)}_n(v)\end{array}$$

Now the induction step has been finished.

Corollary. If $g\in O(F)_n$, then there is a number $k$ such that 
$$g(x_1,\cdots,x_n)\le u^{(k)}_n(\max(x_1,\cdots,x_n))$$

Theorem 3.3. $u_0(x)\in O(F)_0, u_{n+1}(x)\in O(F)_{n+1}\setminus O(F)_n, (n\ge 0).$

{\it Proof}. Firstly we show that $u_{n+1}(x)\in O(F)_{n+1}$.

By induction on $n$. For $n=0$, the result is obvious.

If $u_n(x)\in O(F)_n$, i.e. there exists a $Q\in O(F)$-$LOOP_n$
which can compute $u_n(x)$, then because of 
$u_{n+1}(x)=u^{(x)}_n(1)$, the following program $P$ can compute 
$u_{n+1}(x)$
$$\begin{array}{l}  
Z_1\longleftarrow X_1\\
X_1\longleftarrow 1\\
LOOP\ \ Z_1\\
\ \ Q\\
X_1\longleftarrow Y\\
END\end{array}$$
of course, $P\in O(F)$-$LOOP_{n+1}$ and hence
$u_{n+1}(x)\in O(F)$-$LOOP_{n+1}$.

Now, we show $u_{n+1}(x)\notin O(F)_n$. 

If $u_{n+1}(x)\in O(F)_n$, then there is a $k$ such that
$$u_{n+1}(x)\le u^{(k)}_n(x)$$
but
$$u^{(k)}_n(x)<u_{n+1}(x)\ \ a.e.$$
hence
$$u_{n+1}(x)<u_{n+1}(x)\ \ a.e.$$
a contradiction. So we have proved that $u_{n+1}(x)\notin 
O(F)_n$.

Theorem 3.4. If $P$ is an $O(F)$-$LOOP$ program that can compute
$g(x_1,\cdots,x_m)$ and for some $n\ge 2$, there is a $k$ such 
that
$$ST_P(x_1,\cdots,x_m)\le u^{(k)}_n(\max(x_1,\cdots,x_m)),$$
then $g\in O(F)_n$.

The proof is completely like that of the same result for the
$LOOP$ program in [1]. For details, the reader is refered to [1].

\section{Speed-up Hierarchy and the Speed-up Theorem}

In the rest of this paper, if not stated apparently we will use
$S_i(x)$ to denote the space complexity of Turing machine $M_i$,
and all the results is proved only for space complexity, but we
can easily modify these results to any Blum complexity, especially
to time complexity. for example, the lemma4.1, theorems4.1 and 4.2
are also true to Blum complexity.
 
Definition 4.1. An off-line $TM$ (i.e. Turing Machine) $M$ is a 
multitape $TM$ who has one read-only tape with endmarks and $k$ 
semi-infinite storage tapes. If for every input word of length 
$n$, $M$ scans at most $S(n)$ cells on any storage tape, then $M$ 
is said to be an $S(n)$ space-bounded Turing Machine, or of space 
complexity $S(n)$. The language recognized by $M$ is also said to 
be of space complexity $S(n)$.

Definition 4.2. A function $S(n)$ is said to be space 
constructible if there is some $TM$ $M$ which is $S(n)$ space
bounded and for each $n$, there is some input of length $n$ on 
which $M$ actually uses $S(n)$ tape cells, and especially, if for 
all $n$, $M$ in fact uses exactly $S(n)$ cells on any input of 
length $n$, then we say $S(n)$ is fully space constructible.

Proposition 4.1. For any totally defined recursive function 
$f(x)$, there is a function $g(x)\ge f(x)$ such that $g(x)$ is
space constructible.

Proposition 4.2. $\log n, n^k,2^n$ and $n!$ are all space 
constructible.

Proposition 4.3. If $S_1(n)$ and $S_2(n)$ are space 
constructible, then $S_1(n)\cdot 
S_2(n)$, $2^{S_1(n)}$, $S_1(n)^{S_2(n)}$ and $f(n)=S_1^{(n)}(1)$ are
all space constructible.

Proposition 4.4. If $S(n)\ge n$ is space constructible, then 
$S(n)$ is also fully space constructible.

The proofs of all these Propositions mentioned above are refered 
to [2].

Definition 4.3. The function set sequence ${\cal F}_0,{\cal F}_1,
\cdots$ is said to be a Speed-up hierarchy of function set
${\cal F}$, if it satisfies the following properties (1) to (4)

(1). ${\cal F}=\cup^\infty_{n=0}{\cal F}_n$,

(2). ${\cal F}_0\subset{\cal F}_1\subset\cdots$,

(3). there is a fully space constructible function $u_0(x)$ such
that for any $k$, $u^{(k)}_0(x)\in{\cal F}_0$, $u^{(k)}_n(x)
F_n\setminus{\cal F}_{n-1}$, where $u_n(x)=u^{(x)}_{n-1}(1)$. And 
$\{u_n(x)\}$ are called the hierarchy functions of this Speed-up
hierarchy,

(4). for any $f(x)\in{\cal F}_n$, there is a $k$ such that
$$f(x)\le u^{(k)}_n(x)$$

Example 4.1. The Grzegorczyk hierarchy of the primitive recursive 
functions is a Speed-up hierarchy and the hierarchy functions are 
the functions derived from Ackermann function, i.e. $\{u_n(x)|u_n(x)=A(n,x)\}$.

Example 4.2. Choose a proper function set $F$ in \S2 such that
$$u_0(x)=\max\{f_1(x),f_2(x,x),\cdots,f_m(x,\cdots,x)\}$$
is fully space constructible, then $\{O(F)_n|n=0,1,\cdots\}$ is a 
Speed-up hierarchy of $O(F)$.

Definition 4.4. If $r(x)$ is a totally defined recursive 
function, $L$ is a language on $\sum=\{0,1\}$, then we say that 
$L$ can be speeded up by $r(x)$, if and only if for any $TM$ 
$M_i$ that accepts $L$, there is another $TM$ $M_j$ accepting $L$
such that
$$r(S_j(x))\le S_i(x)$$
where $S_i(x), S_j(x)$ are the space complexity of $M_i$ and 
$M_j$, respectively.

Definition 4.5. We say that language $L$ has ${\cal F}$ complexity 
if there are an $TM$ $M_i$ accepting $L$ and a function 
$f(x)\in{\cal F}$, such that
$$S_i(x)\le f(x)$$

Lemma 4.1. Suppose that $\{u_n(x)\}$ are the hierarchy functions of
some Speed-up hierarchy of ${\cal F}$; for some given $m$, if
$u_m(x)\ge x^2$ a.e., then there is a language $L$ that satisfies 
the following properties (1) to (2).

(1). if $M_i$ accepts $L$, then $S_i(x)>u_{m+1}(x-i)$ a.e.,



(2). for any $k$, there is a $TM$ $M_j$ accepting $L$ such that
for any $x\ge k$
$$S_j(x)\le u_{m+1}(x-k).$$

{\it Proof}. See \S5. Appendix.

Theorem 4.1. If ${\cal F}=\cup^\infty_{n=0}{\cal F}_n$ is a Speed-up 
hierarchy of ${\cal F}$, $\{u_n(x)\}$ are the hierarchy functions
and $u_m(x)\ge x^2$ a.e. for some given $m$, then

(1). there is an $r(x)\in{\cal F}_m$ such that for any nontrivial
language $L$ (i.e. having at least linear complexity) having 
${\cal F}_m$ complexity, $L$ cannot be speeded up by $r(x)$.


(2). for any nondecreasing function $r(x)\in{\cal F}_m$, there is a
language $L$ having ${\cal F}_{m+1}$ complexity such that $L$ can
be speeded up by $r(x)$.



Before proving this theorem, we firstly give some remarks on this 
theorem. Result (1) of this theorem says that Speed-up theorem is 
a hierarchic speed-up: for any speed-up factor $r(x)$, the language
that can be speeded up by $r(x)$ has at least ${\cal F}_{m+1}$ 
complexity, i.e. $L$ has one level higher complexity than $r(x)$.
Result (2) is exactly the Speed-up theorem. Generally speaking, 
results (1) and (2) show that for any $r(x)\in{\cal F}_{m-1}$,
there is a language $L$ having ${\cal F}_m$ complexity which can be
speeded up by $r(x)$. In fact, this speed-up is not a
``real'' speed-up to $L$, becuase the new algorithm obtained from
the speed-up does not have ${\cal F}_{m-1}$ or even lower complexity, but still 
has the ${\cal F}_m$
complexity, i.e. speed-up cannot pull down the complexity in
quantity level. Intuitively, we can say that for any speed-up 
factor $r(x)$, there is a language $L$ whose complexity $S_i(x)$ 
is so high that $r(S_i(x))$ and $S_i(x)$ have the same 
quantitative level of complexity. So it seems that Speed-up 
theorem has not much ``actual'' meaning. For example,useing Grzegorczyk
hierarchy, we can easily show: when $r(x)$ is a polynomial, then
there is a language having exponential complexity that can be speeded
up by $r(x)$, and really the language that can be speeded up
by $r(x)$ has at least exponential complexity and the new
algorithm obtained from the speed-up still
has exponential complexity. Meanwhile, we expect that this 
theorem will be a help to the last solution of $P=?NP$ problem. For
if we can find an $NP$-complete problem which can be speeded up
by a polynomial, then we can assert that $P\ne NP$! The reason
for this is that from the analysis mentioned above we find out that
the language that can be speeded up by a polynomial has no 
algorithm of polynomial complexity(attention: here we have only used
the time complexity and the Grzegorczyk hierarchy).

{\it Proof of theorem} 4.1


(1). Let $r(x)=u_m(x)$, then for any nontrivial language $L$ having
${\cal F}_m$ complexity, there are an $r_m(x)\in{\cal F}_m$ and
an $TM$ $M_i$, such that $M_i$ accepts $L$ and $S_i(x)\le 
r_m(x)$. $r_m(x)\in{\cal F}_m$ shows that there is a number $k$ so 
that $r_m(x)\le u^{(k)}_m(x)$. If $L$ can be speeded up $r(x)$, 
then there is another $TM$ $M_j$ accepting $L$ and 
$r^{(k+1)}(S_j(x))\le S_i(x)\le u^{(k)}_m(x)$, but $S_j(x)\ge x$, 
so $u^{(k+1)}_m(x)\le u^{(k)}_m(x)$, which is a contradiction to 
Proposition 2.6.. Now we've proved that $L$ cannot be speeded up
by $r(x)$.

(2). Proposition 4.1. shows that there is a language $L$ that
satisfies the following properties:

$<$a$>$. If $M_i$ accepts $L$, then $S_i(x)>u_{m+1}(x-i)$ a.e.,
$<$b$>$. For any $k$, there is an $TM$ $M_j$ that accpets $L$ and
$$S_j(x)\le u_{m+1}(x-k);$$


and $r(x)\in{\cal F}_m$ shows that there is a $k_0$ such that
$$r(x)\le u^{(k_0)}_m(x).$$

Now for any $TM$ $M_i$ that accepts $L$, $<$a$>$. shows that
$$S_i(x)>u_{m+1}(x-t)\ \ \ \ a.e.,$$
let $k=i+k_0+1$, then $<$b$>$. shows that there is an $TM$ $M_j$
accpeting $L$ such that
$$S_j(x)\le u_{m+1}(x-k)$$
hence
$$\begin{array}{l}   
r(S_j(x))\le u^{(k_0)}_m(u_{m+1}(x-k))\\
=u_{m+1}(x-k+k_0)\\
<u_{m+1}(x-i)\\
<S_i(x).\end{array}$$

Theorem 4.2. (Internal Displacement Speed-up). Let ${\cal F}
=\cup^\infty_{n=0}{\cal F}_n$ be a Speed-up hierarchy, then for
any $m$ and $c$, there is a language $L$ having ${\cal F}_{m+1}$ 
complexity, such that for any $TM$ $M_i$ accepting $L$, there
is another $TM$ $M_j$ accepting $L$ and
$$S_j(n)\le S_i(n-c)\ \ \ \ a.e.$$

{\it Proof}. Lemma 4.1. shows that for this $m$ there is a language $L$
having ${\cal F}_{m+1}$ complexity and

(1). for any $M_i$ that accepts $L$, we have
$$S_i(n)>u_{m+1}(n-i)\ \ \ \ a.e.,$$

(2). for $k=i+c$, there is an $TM$ $M_j$ that accepts $L$ and
$$S_j(n)<u_{m+1}(n-i-c)$$
so we have
$$S_i(n-c)>u_{m+1}(n-i-c)>S_j(n)\ \ \ \ a.e.$$

When we've obtained the theorems 4.1. and 4.2., perhaps we'll ask: Is
there an internal linear Speed-up theorem?, i.e.

Internal Linear Speed-up Phenomenon: For any $c\le 1,$ there is a language $L$, such that for any $TM$
$M_i$ accepting $L$, there is another $TM$ $M_j$ that
satisfies the property:
$$S_j(n)\le S_i([c*n])\ \ \ \ a.e.$$

About this question and the relations 
among all these Speed-up 
theorems mentioned in this paper, 
we'll deal with in another paper.

\section{Appendix}

{\it Proof of Lemma} 4.1.

Remark: this proof is mainly from [2].
Because $u_0(x)$ is fully space constructible, we can easily show 
that $u_{m+1}(x)$ is also fully space constructible, the details 
is left to reader.

Let $M_1,M_2,\cdots$ be an enumeration of
 all off-line $TM$'s.
Now let us construct $L\subseteq 0^*$ to satisfy (1) and (2). For
$n=0,1,2,\cdots$ in turn, we specify whether $0^n$ is in $L$. In
the process, certain $M_i$ are designated as ``canceled''. A 
canceled $TM$ surely does not accept $L$, let $\sigma(n)$ be the 
least integer $j\le n$ such that $S_j(n)<u_{m+1}(n-j)$, and $M_j$ 
is not canceled by $i=0,1,\cdots,n-1$. When we consider $n$, if 
$\sigma(n)$ exists, $M_{\sigma(n)}$ is designated as canceled. 
Then $0^n$ is placed in $L$ if and only if $\sigma(n)$ exists and
$0^n$ is not accepted by $M_{\sigma(n)}$.

Next we prove that $L$ satisfies (1), namely: if $L(M_i)=L$, then 
$S_i(n)\ge u_{m+1}(n-i)$ a.e. Let $L(M_i)=L$. In constructing 
$L$, all $TM$'s $M_j$, for $j<i$, that are ever canceled are 
canceled after considering some finite number of $n$'s, say up to 
$n_0$. Note that $n_0$ cannot be effectively computed, but 
nevertheless exists. Suppose $S_i(n)<u_{m+1}(n-i)$ for some 
$n>\max(n_0,i)$. When we consider $n$, no $M_j, j<i$ is canceled. 
But an $TM$ that is canceled will surelly not accept $L$. Thus 
$S_i(n)\ge u_{m+1}(n-i)$ for $n>\max(n_0,i)$, that is, $S_i(n)\ge 
u_{m+1}(n-i)$ a.e..

To prove condition (2), we show that there exists, for given $k$, 
an $TM$ $M=M_j$ such that $L(M)=L$, and $S_j(n)\le u_{m+1}(n-k)$
for all $n$. To determine whether $0^n$ in $L$, $M$ must simulate
$M_{\sigma(n)}$ on $0^n$. To know what $\sigma(n)$ is, $M$ must
determine which $M_i$'s have already been canceled by $0^l$ for
$l<n$. However, constructing the list of canceled $TM$'s directly 
requires seeing if $M_i$ uses more than $u_{m+1}(l-i)$ space for 
$0\le l\le n$ and $1\le i\le n$, this requires more than 
$u_{m+1}(n-k)$ space.

The solution is to observe that any $TM$ $M_i$, $i\le k$, that is 
ever canceled, is canceled when we consider some $l$ less than a 
particular $n_1$. For each $l\le n_1$, incorporate into the 
finite control of $M$ whether $0^l$ is in $L$, and also
incorporate a list of all $TM$'s $M_i$ canceled by $l\le n_1$. 
Thus no space at all is needed by $M$ if $n\le n_1$. If $n>n_1$, 
to compute $\sigma(n)$ and simulate $M_{\sigma(n)}$, on $0^n$, it
will only be necessary to simulate $TM$'s $M_i$ on input $0^l$,
where $n_1<l\le n$ and $k<i\le n$, to see whether $M_i$ is 
canceled by $l$.

To test whether $M_i$ is canceled by $l$, we need only simulate 
$M_i$ using $u_{m+1}(l-i)$ of $M_i$'s cells, which is less than 
$u_{m+1}(n-k)$, as $l\le n$ and $i>k$. As $n>n_1$, it must be 
that $\sigma(n)$, if it exists, is greater than $k$. Thus 
simulating $M_{\sigma(n)}$ on input $0^n$ takes $u_{m+1}(n-
\sigma(n))$ of $M_{\sigma(n)}$'s cells, which is less than 
$u_{m+1}(n-k)$ cells.

Lastly, we must show that $M$ can be made to operate within space 
$u_{m+1}(n-k)$. We need only simulate $TM$'s $M_i$ for $k<i\le n$ 
on inputs $0^l$, $n_1\le l\le n$, to see whether they get
canceled, so we need represent no more than $u_{m+1}(n-k-1)$ 
cells of $M_i$'s tape for any simulation. Since $i\le n$, the 
integer code for $M_i$ has length no more than $\log_2 n$. Thus
any tape symbol of $M_i$ can be coded using $\log_2 n$ of $M$'s
cells. As
$$u_{m+1}(n-k)\ge [u_{m+1}(n-k-1)]^2\ge\log_2 n\cdot u_{m+1}(n-k-
1)\ \ a.e.,$$
$u_{m+1}(n-k)$ space is sufficient for the simulation for almost 
all $n$.

In addition to the sapce required for simulating the $TM$'s, 
space is needed to maintain the list of canceled $TM$'s. This 
list consists of at most $n$ $TM$'s, each with a code of length 
at most $\log_2 n$. The $n\cdot\log_2 n$ space needed to maintain
the list of canceled $TM$'s is also less than $u_{m+1}(n-k)$ 
a.e., so $M$ can be modified to recognize words $0^n$, where
$u_{m+1}(n-k)<[u_{m+1}(n-k-1)]^2$, or $u_{m+1}(n-k-1)<\log_2 n$, 
or $n\cdot\log_2 n>u_{m+1}(n-k)$, in its finite control. The 
resulting $TM$ is of sapce complexity $u_{m+1}(n-k)$ for all $n$, 
and is the desired $M_j$.

\section{Acknowledgement}

Before finishing this paper, we must 
mention that professor Xu 
Shuruen advised us to give more attention on 
the reason why there 
is the Speed-up phenomenon. And this paper is the result of
repeadedly discussions between professor Xu and the author. Here 
the author would like to express his great 
gratitute to professor Xu.

\begin{thebibliography}{99}

\bibitem{[1]} 
M.D.Davis, E.J.Weyuker, Computability, Complexity, and 
Languages, Academic Press. 1983.

\bibitem{[2]} J.E.Hopcroft, J.D.Ullman, 
Introduction to Automata Theory, 
Languages, and computation, 1979.

\bibitem{[3]}Cutland, Computability, 1981.

\bibitem{[4]}  E.Borger, 
Computability, Complexity, Logic, North-Horland. 
1989.


\bibitem{[5]} E.Rose, Subrecursion: 
Functions and Hierarchy, Oxford Logic 
Series. 1984.

\bibitem{[6]} Grzegorczyk, Some classes of 
recursive functions. Selected 
Papers on Mathematical Logic.

\bibitem{[7]} Zhang Min-hua, 
Computability, Qin Hua University
Press. 1984.
\end{thebibliography}
\end{document}


\bibitem{[1]} 




